\chapter{Predictive Benchmarks}
\label{ch:ml_extension}

\section{Purpose}

This appendix benchmarks one-year-ahead prediction using lagged country-level predictors. The task is intentionally modest because the panel is small and country structure is persistent.

\section{Persistence Baseline}

The decisive benchmark is simple persistence. Table~\ref{tab:ml_naive_baselines} shows that previous-year unmet need has the lowest held-out test MAE among the naive and statistical baselines and beats the flexible models in Table~\ref{tab:ml_model_performance_full}.

\begin{table}[H]
\centering
\small
\caption{Naive and statistical baselines for one-year-ahead prediction. The held-out test set is the time-ordered 2021--2023 complete-case prediction sample; predictors use prior-year country information only.}
\label{tab:ml_naive_baselines}
\begin{tabular}{@{}rrrr@{}}
\toprule
Model & MAE & RMSE & R2 \\
\midrule
Previous-year outcome & 0.587 & 1.006 & 0.850 \\
AR(1) linear & 0.661 & 1.101 & 0.820 \\
Ridge with lagged predictors & 0.695 & 1.115 & 0.816 \\
Last observed country value & 0.718 & 1.187 & 0.791 \\
Rolling country mean & 0.796 & 1.357 & 0.727 \\
Country mean & 0.855 & 1.340 & 0.734 \\
Last training-year mean & 1.815 & 2.613 & -0.013 \\
Global mean & 1.885 & 2.596 & -0.000 \\
\bottomrule
\end{tabular}

\end{table}

\begin{table}[H]
\centering
\small
\caption{Linear and flexible benchmark performance with lagged predictors. Models are trained on the pre-test complete-case prediction panel and evaluated on the same 2021--2023 held-out test years as the naive baselines. Additional model-inspection tools are not used because the chapter is a negative benchmark, not an algorithmic contribution.}
\label{tab:ml_model_performance_full}
\begin{tabular}{@{}rrrrrrr@{}}
\toprule
Model & Family & CV MAE & CV SD & Test MAE & RMSE & R2 \\
\midrule
Lasso & linear & 0.561 & 0.202 & 0.667 & 1.116 & 0.815 \\
Elastic Net & linear & 0.556 & 0.169 & 0.675 & 1.133 & 0.809 \\
Ridge & linear & 0.572 & 0.159 & 0.693 & 1.139 & 0.807 \\
Random Forest & tree & 0.544 & 0.100 & 0.727 & 1.313 & 0.744 \\
XGBoost & tree & 0.624 & 0.180 & 0.769 & 1.350 & 0.730 \\
\bottomrule
\end{tabular}

\end{table}

\section{Validation Checks}

The evaluation uses temporal train/test separation, expanding-window validation, bootstrap intervals for held-out MAE, and blocked country-holdout validation. Country-holdout validation is especially important because aggregate predictors can encode persistent country identity rather than transferable structure.

\begin{table}[H]
\centering
\small
\caption{Blocked country-holdout predictive performance. Folds hold out countries rather than years to test transfer beyond persistent country structure; no population weights are applied.}
\label{tab:ml_country_holdout_summary}
\begin{tabular}{@{}rrrrr@{}}
\toprule
model & mean\_mae & sd\_mae & mean\_rmse & folds \\
\midrule
Lasso & 0.594 & 0.210 & 0.944 & 5 \\
Elastic Net & 0.621 & 0.202 & 0.969 & 5 \\
Ridge & 0.648 & 0.206 & 0.993 & 5 \\
Random Forest & 0.747 & 0.393 & 1.228 & 5 \\
XGBoost & 0.783 & 0.330 & 1.298 & 5 \\
\bottomrule
\end{tabular}

\end{table}

\section{Interpretation}

Open Eurostat indicators contain next-year predictive signal, but simple persistence is the strongest benchmark in the held-out test. Flexible machine learning therefore remains appendix-level validation evidence rather than an independent thesis contribution.
